% Exact LaTeX recreation of the supplied Phase-1.pdf
% Compile with pdfLaTeX. Place member pictures in the same folder if required.

\documentclass[11pt,a4paper]{article}

\usepackage[a4paper,left=2.0cm,right=2.0cm,top=1.55cm,bottom=1.55cm]{geometry}
\usepackage[T1]{fontenc}
\usepackage{lmodern}
\usepackage{graphicx}
\usepackage{xcolor}
\usepackage{array}
\usepackage{tabularx}
\usepackage{ragged2e}
\usepackage{enumitem}
\usepackage{fancyhdr}
\usepackage{tikz}
\usepackage{setspace}

% -------------------------------------------------
% COLOURS
% -------------------------------------------------
\definecolor{TitleBlue}{RGB}{61,103,154}

% -------------------------------------------------
% GENERAL SETTINGS
% -------------------------------------------------
\setlength{\parindent}{0pt}
\setlength{\parskip}{0pt}
\renewcommand{\arraystretch}{1.0}

% The supplied document uses a clean sans-serif heading style
% and italic serif-style explanatory text.
\newcommand{\blueheading}[1]{%
    {\sffamily\color{TitleBlue}\fontsize{15}{18}\selectfont #1}%
}

\newcommand{\bluebigheading}[1]{%
    {\sffamily\bfseries\color{TitleBlue}\fontsize{16}{19}\selectfont #1}%
}

\newcommand{\instruction}[1]{%
    {\itshape\fontsize{10.5}{12.5}\selectfont #1}%
}

\newcommand{\answerbox}[2][3.0cm]{%
    \vspace{2pt}
    \noindent
    \fbox{%
        \begin{minipage}[t][#1][t]{\dimexpr\textwidth-2\fboxsep-2\fboxrule\relax}
            \vspace{1mm}
            \itshape #2
        \end{minipage}
    }
}

% -------------------------------------------------
% HEADER / FOOTER
% -------------------------------------------------
\pagestyle{fancy}
\fancyhf{}
\renewcommand{\headrulewidth}{0pt}
\renewcommand{\footrulewidth}{0pt}

\fancyhead[L]{%
    \sffamily\bfseries\itshape\fontsize{10.5}{12}\selectfont
    COURSE: B.Tech CSE\\
    TEAM ID: T323
}
\fancyhead[R]{%
    \sffamily\bfseries\fontsize{10.5}{12}\selectfont
    PROJECT PROPOSAL
}
\fancyfoot[R]{%
    \itshape\fontsize{10}{12}\selectfont
    Page \thepage
}
\setlength{\headheight}{14pt}
\setlength{\headsep}{23pt}
\setlength{\footskip}{24pt}

% -------------------------------------------------
% PAGE 1
% -------------------------------------------------
\begin{document}

\vspace*{4mm}

\bluebigheading{PROJECT AND TEAM INFORMATION}

\vspace{13mm}

\blueheading{Project Title}

\vspace{1mm}

\answerbox[1.0cm]{SNAKE GAME LOGIC SYSTEM}

\vspace{13mm}

\blueheading{Student / Team Information}

\vspace{2mm}

% Team information table
% Compact 4-member layout so all members remain on Page 1.
\noindent
\begin{tabularx}{\textwidth}{|>{\RaggedRight\arraybackslash}p{0.69\textwidth}|>{\RaggedRight\arraybackslash}X|}
\hline
\begin{minipage}[t]{\linewidth}
\vspace{1.5mm}
\textit{Team Name:}\\[-1mm]
\vspace{1.5mm}
\end{minipage}
&
\begin{minipage}[t]{\linewidth}
\vspace{1.5mm}
\textit{DSCPP-III-2026-T323}
\vspace{1.5mm}
\end{minipage}
\\
\hline
\begin{minipage}[t][3.25cm][t]{\linewidth}
\vspace{1.5mm}
\textbf{Team member 1 (Team Lead)}\\[-1mm]
\textit{(Anna, Mayank: 2510370093: 2510370093@geu.ac.in, picture):}
\vspace{2.0cm}
\end{minipage}
&
\begin{minipage}[t][3.25cm][t]{\linewidth}
\vspace{1.5mm}
\textit{Anna, Mayank -- 2510370093}\\
\textit{2510370093@geu.ac.in}

\vspace{2mm}
\includegraphics[width=3.0cm,height=1.5cm,keepaspectratio]{}
\end{minipage}
\\
\hline
\begin{minipage}[t][3.25cm][t]{\linewidth}
\vspace{1.5mm}
\textbf{Team member 2}\\[-1mm]
\textit{(Lingwal, Mamta: 2510380154: 2510380154@geu.ac.in, picture):}
\vspace{2.0cm}
\end{minipage}
&
\begin{minipage}[t][3.25cm][t]{\linewidth}
\vspace{1.5mm}
\textit{Lingwal, Mamta -- 2510380154}\\
\textit{2510380154@geu.ac.in}

\vspace{2mm}
\includegraphics[width=3.0cm,height=1.5cm,keepaspectratio]{}

\end{minipage}
\\
\hline
\begin{minipage}[t][3.25cm][t]{\linewidth}
\vspace{1.5mm}
\textbf{Team member 3}\\[-1mm]

\vspace{2.0cm}
\end{minipage}
&
\begin{minipage}[t][3.25cm][t]{\linewidth}
\vspace{1.5mm}
\textit{Singh, Abhinav -- Student ID}\\
\textit{email@example.com}

\vspace{2mm}

\end{minipage}
\\
\hline
\end{tabularx}

\newpage

% -------------------------------------------------
% PAGE 2
% -------------------------------------------------

\bluebigheading{PROPOSAL DESCRIPTION (10 pts)}

\vspace{12mm}

\blueheading{Motivation (1 pt)}

\vspace{5mm}


\answerbox[7.0cm]{Snake game is a simple game, however creation of such a game gives us a practical example on how we can use Data Structures and Algorithms in solving real-world problems. The problem this project needs to solve is how to deal with the continuous movement and size change of snake while solving the collision detection, food generation, score calculation and updating game status. The snake's head should move and its tail should disappear unless it eats the food. The snake should be able to grow its body when it eats the food. And also the game should be able to detect the position of snake so as to avoid any collision and generate food. The importance of this project is that it connects theoretical ideas of DSA to a practical and graphical implementation. For example, theoretical concepts such as linked lists, queues, searching and collision detection are more easily understood when implemented in a practical situation. Furthermore, this work improves the ability to think logically and solve problems as well as to understand how algorithms affect the performance of the program.}
\vspace{12mm}

\blueheading{State of the Art / Current solution (1 pt)}

\vspace{5mm}

\answerbox[7.65cm]{Currently, Snake Game apps are created through coding languages such as C++, Java, Python, or JavaScript, which have data structures and game libraries incorporated into them. The body of the snake is represented using arrays, lists, queues or linked lists where each element represents the location of a body segment of the snake. The position of the head is calculated at each move. The snake is cut off at the tail so that it preserves the same size when it does not eat any food. When the snake eats a piece of food the tail stays and grows the snake. Collisions are detected when the head's position collides with either the boundary of the game or any body part of the snake. Food is randomly generated, but not on the snake. This is done through a game loop that processes input, movement, collisions, rendering and scoring. The modern implementation would involve the use of graphical libraries to make animations and improved user interaction. Despite the fact that current solutions offer fluid gameplay, the data structures utilized behind them might be hidden within library files or high-level functions. This project will aim at implementing all these processes directly so as to illustrate the working of DSA techniques in practical real-time applications.}
\vspace{12mm}

\newpage

% -------------------------------------------------
% PAGE 3
% -------------------------------------------------
\vspace{12mm}

\blueheading{Project Goals and Milestones (2 pts)}

\vspace{5mm}

\answerbox[17.15cm]{General goals and Milestones\\
The main goal of this project is the implementation of an interactive game named Snake game using C++. Data Structures and Algorithms would be applied for efficient handling of data structures.\\

The goal of this project are as follows:\\

-	To design a user-friendly Snake game.\\
-	To implement the snake’s body using appropriate data structure like linked list or queue.\\
-	To implement functions for adding the head, removing the tail, and checking the snake’s position.\\
-	To implement movement and control via keyboard.\\
-	To implement collision with walls and the snake’s body.\\
-	To randomly generate food at valid locations.\\
-	To keep track of the player’s score.\\
-	To implement game-over feature.\\
-	To test the game’s functionality.\\

Project Milestones\\

Milestone 1 - Planning and Requirement Analysis: Determine the needs, rules, functionalities, and DSA concepts of the game.\\

Milestone 2 - Data Structure Design: Design the snake data structure and develop some basic functions like "addHead()", "removeTail()", and "isSnakePosition()".\\

Milestone 3 - Main Development of the Game: Develop the functions of "setup()", "input()", "logic()", and "draw()".\\

Milestone 4 - Features of the Game: Introduce features like adding food, updating score, collision, game-over, and re-initiation.\\

Milestone 5 - Testing and Debugging: Test various kinds of movements, foods, boundaries, self-collision, and restarting of the game.\\

Milestone 6 - Finishing Touches: Enhance the interface of the game, debug all the remaining errors, document the project, and prepare the demonstration and presentation of the final product.\\
}

\newpage

\vspace{12mm}

\blueheading{Project Approach (3 pts)}

\vspace{5mm}

\answerbox[14.05cm]{Solution Design and Technologies\\

The suggested solution would be to create a Snake game in C++ that emphasizes practical use of data structures and algorithms. This game will be designed in such a way that separate functions would take care of the following processes: initial setting of the game, user input, movement of the snake, collision detection, food creation, scoring, and game restart. The snake can be implemented as a linked list/queue where each node/element contains the position of a part of the snake. addHead() function would be used for adding the head of the snake and removeTail() for removal of the tail in case the snake moves without eating food. isSnakePosition() function would help to check if a certain position belongs to the snake. The primary loop of the program will execute three tasks over and over again: input, logic and rendering. Keyboard input will control the direction of the snake. The logic module will compute the next position, detect if the snake has consumed the food, grow the snake, add to the score, and detect collision. The rendering module will show the position of the snake, the food, the score, and the boundaries of the game.\\

Platforms and Technologies Programming Language:\\

C++ Platform: Windows/Linux Computer\\
 Development Environment: Visual Studio Code / Code::Blocks\\
Compiler: GCC (g++)\\
Data Structures: Linked List, Queue \\
Algorithms: Searching, position detection, collision detection, generating random food \\
User Interface: console/terminal user interface\\
The solution will be built gradually, and then tested and debugged. The different conditions within the game will be tested to validate movement, collision detection, snake growing, score and restarting the game.}


\newpage

\vspace{12mm}

\blueheading{System Architecture (High Level Diagram)(2 pts)}

\vspace{1mm}

\noindent
\fbox{%
    \begin{minipage}{\dimexpr\linewidth-2\fboxsep-2\fboxrule\relax}
        \centering
        \includegraphics[
            width=0.95\linewidth,
            height=7cm,
            keepaspectratio
        ]{WhatsApp Image 2026-09-02 at 8.36.27 AM.jpeg}
    \end{minipage}%
}

\vspace{5mm}
  

\vspace{5mm}


\blueheading{Project Outcome / Deliverables (1 pts)}

\vspace{5mm}

\answerbox[10.7cm]{Outcomes of the Project Deliverables\\

The main outcome of this project is going to be the completion of the Snake Game that uses Data Structures and Algorithms in C++.\\

The project will provide the following outcomes:\\
- A working Console-based Snake Game where the snake moves smoothly and can be controlled by the user. \\
- The implementation of the snake using the appropriate data structure, e.g., linked list or queue. \\
- The use of functions like addHead(), removeTail() and isSnakePosition() in managing the snake. \\
- Implementation of all the basic functionalities of the game, e.g., food creation, snake growing, score, etc. \\
- Game over and restarting mechanisms. - Simple but effective display of the game on the console including the snake, food, boundary, score, etc. \\
- Effective management of various game states, e.g., collision with walls and itself, etc. - Working source code of the game implemented in C++. \\
- Documentation of the project which covers the problem statement, objectives, data structures, al-gorithms, functions, and working of the system. \\
- Demonstration of the completed project/presentation.}

\newpage

\vspace{12mm}

\blueheading{Assumptions}

\vspace{5mm}

\answerbox[6.65cm]{Assumptions\\
These assumptions are made when designing the Snake Game:\\
 - The game is played by only one person.\\
- The game is designed to run in a computer with a C++ compiler. \\
- Keyboard inputs are available to control the snake. \\
- The game board has predefined dimensions and boundaries. \\
- The snake moves in steps with a fixed speed. \\
- Food appears at random places inside the boundaries of the game. \\
- Food does not appear on the body of the snake. \\
- The player loses the game when the snake touches the boundary or its own body. \\
- Score increases whenever the snake eats the food. \\
- Option to restart the game after game over is provided.}

\vspace{12mm}

\blueheading{References}

\vspace{5mm}

\answerbox[6.65cm]{Sources / References\\

GeeksforGeeks – Snake Game in C \\
For understanding the basics of how a Snake Game works in the console window, like movements, generating food, score, and detecting collisions. \\
geeksforgeeks.org\\

cppreference – std::queue\\
 For understanding how to perform queue operations like push(), pop(), front(), and back(). cppreference – cppreference.com \\

cppreference – std::list \\
For understanding linked list and its operations like insertions and removals.}

\end{document}
